\subsection{在模理论中的应用}\label{src:A.6.2-subsection}
下面是上同调与基变换在模理论中的一个典型应用, 它建立光滑曲线族的性质 (命题 \bookref{6.1.16}), 例如用于 $\mathcal M_g$ 的代数性 (\bookref{4.1.17}). 以下论证同样适用于稳定曲线族 (命题 \bookref{6.3.23}).
\begin{proposition}\label{src:A.6.11}
设 $\pi:C\to S$ 为亏格 $g\geq2$ 的光滑曲线族 (即 $C\to S$ 为光滑固有概形态射, 每条几何纤维都是亏格 $g$ 的连通曲线). 则
\begin{enumerate}[(1)]
\item $\pi_*\mathcal O_C=\mathcal O_S$;
\item $\pi_*(\Omega_{C/S})$ 为秩 $g$ 的向量丛, 构造与 $S$ 上的基变换交换, 且 $\mathrm R^1\pi_*(\Omega_{C/S})\simeq\mathcal O_S$, 而对 $i\geq2$ 有 $\mathrm R^i\pi_*(\Omega_{C/S})=0$;
\item 对 $k>1$, 推前 $\pi_*(\Omega_{C/S}^{\otimes k})$ 为秩 $(2k-1)(g-1)$ 的向量丛, 构造与 $S$ 上的基变换交换, 且对 $i>0$ 有 $\mathrm R^i\pi_*(\Omega_{C/S}^{\otimes k})=0$.
\end{enumerate}
\end{proposition}
\begin{proof}
对于 (1), 注意对所有 $s\in S$, $\mathrm H^0(C_s,\mathcal O_{C_s})=\kappa(s)$, 因为 $C_s$ 固有且几何连通. 因而 $\phi_s^0:\pi_*\mathcal O_C\otimes\kappa(s)\to\mathrm H^0(C_s,\mathcal O_{C_s})$ 满射. 在上同调与基变换 (\bookref{A.6.8}(a)--(b)) 中取 $i=0$, 得到 $\phi_s^0$ 为同构, 且 $\pi_*\mathcal O_C$ 为线丛. 在 $s\in S$ 上的纤维上, 自然映射 $\mathcal O_S\to\pi_*\mathcal O_C$ 诱导满射 $\kappa(s)\to\pi_*\mathcal O_C\otimes\kappa(s)$ (因为后接 $\phi_s^0:\pi_*\mathcal O_C\otimes\kappa(s)\to\mathrm H^0(C_s,\mathcal O_{C_s})=\kappa(s)$ 得到恒等映射). 因而 $\mathcal O_S\to\pi_*\mathcal O_C$ 为线丛之间的满态射, 从而为同构.

对于 (2), 由于 $\Omega_{C/S}$ 为相对对偶化层 (见 \cite[\S$6.4$]{Liu02}), Grothendieck--Serre 对偶性意味着 $\mathrm R^1\pi_*\Omega_{C/S}\simeq\pi_*\mathcal O_C$, 由 (1) 将其与 $\mathcal O_S$ 等同. 对 $i\geq2$, 有 $\mathrm H^i(C_s,\Omega_{C/S}\otimes\kappa(s))=0$ (因为 $\dim C_s=1$), 而 \bookref{A.6.8}(a) 意味着 $\mathrm R^i\pi_*\Omega_{C/S}=0$. 对 $i=2$ 应用 \bookref{A.6.8}(b), 得到对每个 $s\in S$, $\phi_s^1:\mathrm R^1\pi_*\Omega_{C/S}\otimes\kappa(s)\to\mathrm H^1(C_s,\Omega_{C_s/\kappa(s)})$ 满射; 再对 $i=1$ 应用 \bookref{A.6.8}(a), 可知 $\phi_s^1$ 为同构. 由于 $\mathrm R^1\pi_*\Omega_{C/S}$ 为线丛, 对 $i=1$ 应用 \bookref{A.6.8}(b), 可知 $\phi_s^0$ 满射; 再对 $i=0$ 应用 \bookref{A.6.8}(a)--(b), 得到 $\pi_*\Omega_{C/S}$ 为秩 $\mathrm h^0(C_s,\Omega_{C_s/\kappa(s)})=g$ 的向量丛, 构造与基变换交换.

对于 (3), 由于 $k>1$, 对每个 $s\in S$ 有 $\deg(\Omega_{C_s/\kappa(s)}^{\otimes(1-k)})<0$, 而 Serre 对偶性 (\bookref{6.1.3}) 意味着 $\mathrm H^1(C_s,\Omega_{C_s/\kappa(s)}^{\otimes k})=\mathrm H^0(C_s,\Omega_{C_s/\kappa(s)}^{\otimes(1-k)})=0$. 注意对 $i\geq2$, $\mathrm H^i(C_s,\Omega_{C_s/\kappa(s)}^{\otimes k})=0$, 因为 $\dim C_s=1$. 上同调与基变换 (\bookref{A.6.8}(a)) 给出对 $i>0$, $\mathrm R^i\pi_*(\Omega_{C/S}^{\otimes k})=0$. 另一方面, 由 Riemann--Roch (\bookref{6.1.2}), 有 $\mathrm h^0(C_s,\Omega_{C_s/\kappa(s)}^{\otimes k})=\deg(\Omega_{C_s/\kappa(s)}^{\otimes k})+1-g=(2k-1)(g-1)$. 对 $i=0$ 应用上同调与基变换 (\bookref{A.6.8}(b)), 得到 $\pi_*(\Omega_{C/S}^{\otimes k})$ 为秩 $(2k-1)(g-1)$ 的向量丛.
\end{proof}
类似地, 可用上同调与基变换建立凝聚层族的性质, 例如在 $\cBun(C)$ 的代数性 (\bookref{4.1.23}) 中需要它.
\begin{proposition}\label{src:A.6.12}
设 $p:X\to S$ 为固有概形态射, $F$ 为 $X$ 上在 $S$ 上平坦的有限表现拟凝聚层. 假设对所有 $s\in S$ 有 $\dim X_s\leq d$. 满足对所有 $j>0$ 有 $\mathrm H^j(X_s,F_s)=0$ 的点 $s\in S$ 组成的子集 $S'$ 为开集. 记 $X'=p^{-1}(S')$, $p':=p|_{X'}:X'\to S$, $F'=F|_{X'}$, 则对所有 $j>0$ 有 $\mathrm R^jp'_*F'=0$, 且 $p'_*F'$ 为构造与基变换交换的向量丛.
\end{proposition}
\begin{proof}
对每个 $j=1,\ldots,d$, \bookref{A.6.8}(a) 意味着满足 $\mathrm H^j(X_s,F_s)=0$ 的点 $s\in S$ 组成开集, 且比较映射 $\phi_s^j:\mathrm R^jp_*F\otimes\kappa(s)\to\mathrm H^j(X_s,F_s)$ 为同构. 因而 $\mathrm R^jp'_*F=0$, 可以对 $i=1$ 应用 \bookref{A.6.8}(b), 得到 $\phi_s^0:p'_*F'\otimes\kappa(s)\to\mathrm H^0(X_s,F_s)$ 满射. 对 $i=0$ 应用 \bookref{A.6.8}(a)--(b), 即得最后的结论.
\end{proof}

\subsection{在线丛中的应用}\label{src:A.6.3-subsection}
给定固有平坦态射 $f:X\to Y$, $X$ 上的线丛 $L$ 何时是 $Y$ 上某个线丛的拉回? 更一般地, 是否存在最大的子概形 $Z\subseteq Y$, 使得 $X_Z=X\times_Y Z$ 上的 $L_Z$ 为 $Z$ 上某个线丛的拉回? 本节按复杂程度递增的顺序给出三个答案. 由于这些结果依赖纤维 $X_y$ 的性质, 先讨论各种条件之间的关系.
\begin{lemma}\label{src:A.6.13}
设 $f:X\to Y$ 为 Noether 概形之间的固有平坦态射. 考虑以下条件:
\begin{enumerate}[(1)]
\item $f:X\to Y$ 的几何纤维非空, 连通且约化;
\item 对所有 $y\in Y$, $\mathrm h^0(X_y,\mathcal O_{X_y})=1$;
\item $\mathcal O_Y=f_*\mathcal O_X$, 且在任意基变换后仍成立 (即对概形态射 $T\to Y$, 有 $\mathcal O_T=f_{T,*}\mathcal O_{X_T}$).
\end{enumerate}
则 (1) $\Rightarrow$ (2) $\Longleftrightarrow$ (3).
\end{lemma}
\begin{proof}
若 (1) 成立, 则由平坦基变换 (\bookref{A.2.12}), 有
\[
\mathrm H^0(X_y,\mathcal O_{X_y})\otimes_{\kappa(y)}\overline{\kappa(y)}=\mathrm H^0(X\times_Y\overline{\kappa(y)},\mathcal O_{X\times_Y\overline{\kappa(y)}}).
\]
由于代数闭域上连通, 约化且固有的概形只有常值函数, 得到 $\mathrm h^0(X_y,\mathcal O_{X_y})=1$. 若 (2) 成立, 则比较映射 $\phi_y^0:f_*\mathcal O_X\otimes\kappa(y)\to\mathrm H^0(X_y,\mathcal O_{X_y})=\kappa(y)$ 必为满射, 因为有全局截面 $1\in\mathrm H^0(Y,f_*\mathcal O_X)$. 对 $i=0$ 应用定理 \bookref{A.6.8}, 得到 $f_*\mathcal O_X$ 为线丛. 由于 $\mathcal O_Y\to f_*\mathcal O_X$ 为线丛之间的满射, 它是同构. 同样的论证适用于基变换 $X_T\to T$, 故得 (3). 反之, 考虑映射 $T=\Spec\kappa(y)\to Y$, 即得 (3) $\Rightarrow$ (2).
\end{proof}
当基底约化时, Grauert 定理完整地回答了线丛何时为拉回的问题.
\begin{proposition}[版本 $1$]\label{src:A.6.14}
设 $f:X\to Y$ 为 Noether 概形之间的固有平坦态射, 对所有 $y\in Y$ 满足 $\mathrm h^0(X_y,\mathcal O_{X_y})=1$. 设 $L$ 为 $X$ 上的线丛. 若 $Y$ 约化, 则存在 $Y$ 上的线丛 $M$ 满足 $L=f^*M$, 当且仅当对所有 $y\in Y$, $L_y$ 平凡. 此外, 若这些条件成立, 则 $M=f_*L$, 且伴随态射 $f^*f_*L\to L$ 为同构.
\end{proposition}
\begin{proof}
几何纤维上的条件意味着 $\mathrm h^0(X_y,L_y)=1$, 而 Grauert 定理 (\bookref{A.6.7}) 意味着 $f_*L$ 为线丛, 且 $f_*L\otimes\kappa(y)\xrightarrow{\sim}\mathrm H^0(X_y,L_y)$ 为同构. 断言 $f^*f_*L\to L$ 满射. 只需证明 $(f^*f_*L)|_{X_y}\to L|_{X_y}$ 满射. 记 $f_y:X_y\to\Spec\kappa(y)$, 则有等同
\[
(f^*f_*L)|_{X_y}=f_y^*(f_*L\otimes\kappa(y))=f_y^*\mathcal O_{\Spec\kappa(y)}=\mathcal O_{X_y},
\]
从而断言成立. 由于 $f^*f_*L\to L$ 为线丛之间的满射, 它是同构.
\end{proof}
\begin{exercise}\label{src:A.6.15}
证明若 $Y$ 为连通且约化的 Noether 概形, $E$ 为 $Y$ 上的向量丛, 则 $\operatorname{Pic}(\mathbb P(E))=\operatorname{Pic}(Y)\times\mathbb Z$. 另见 \cite[习题 III.$12.5$]{Har77}.
\end{exercise}
\begin{proposition}[版本 $2$]\label{src:A.6.16}
设 $f:X\to Y$ 为具有整几何纤维的 Noether 概形之间的固有平坦态射. 对 $X$ 上的线丛 $L$, 点集 $\{y\in Y\mid L_y\text{ 在 }X_y\text{ 上平凡}\}$ 为 $Y$ 的闭子集.
\end{proposition}
\begin{proof}
关键观察是, 对域 $\Bbbk$ 上几何整且固有的概形 $Z$, 线丛 $M$ 平凡, 当且仅当 $\mathrm h^0(Z,M)>0$ 且 $\mathrm h^0(Z,M^\vee)>0$. 为证明后一条件充分, 注意有非零同态 $\mathcal O_Z\to M$ 和 $\mathcal O_Z\to M^\vee$, 后者对偶化为非零映射 $M\to\mathcal O_Z$. 由于 $Z$ 为整概形, 复合 $\mathcal O_Z\to M\to\mathcal O_Z$ 也非零, 因而为同构, 因为它对应于 $\mathrm H^0(Z,\mathcal O_Z)=\Bbbk$ 中的非零常数. 于是 $M\to\mathcal O_Z$ 为线丛之间的满射, 从而为同构. 由半连续性定理 (\bookref{A.6.4}), 条件 $\mathrm h^0(X_y,L_y)>0$ 和 $\mathrm h^0(X_y,L_y^\vee)>0$ 是闭条件, 即得结论. 另见 \cite[推论 $6$, 第 $54$ 页]{Mum70a}.
\end{proof}
\begin{remark}\label{src:A.6.17}
若几何纤维仅连通且约化, 则上述点集未必闭. 例如, 给定光滑曲线上的光滑曲线族 $f:X\to Y$, 考虑在闭点 $x\in X$ 处的爆破 $\operatorname{Bl}_xX\to X$, 例外除子为 $E$. 则 $\operatorname{Bl}_xX\to Y$ 为固有平坦态射, 而 $f(x)\in Y$ 上的纤维连通且约化, 但可约. 令 $L=\mathcal O_{\operatorname{Bl}_xX}(E)$, 则纤维 $L_y$ 平凡当且仅当 $y\neq f(x)$.
\end{remark}
对模理论的应用, 必须允许基底 $Y$ 非约化, 并给点集 $Z\subseteq Y$ 一个函子描述. 许多应用 (例如稳定曲线族) 还需要允许可约纤维 $X_y$. 最后这个最强的版本综合了上面两个版本, 并用上同调与基变换的代数表述 (\bookref{A.6.2}) 证明. 此结果将用于 $\mathcal M_g$ 的代数性 (\bookref{4.1.17}) 的证明, 以给出 Hilbert 概形中参数化三重典范嵌入的光滑曲线的局部闭子概形.
\begin{proposition}[版本 $3$]\label{src:A.6.18}
设 $f:X\to Y$ 为 Noether 概形之间的固有平坦态射, 对所有 $y\in Y$ 满足 $\mathrm h^0(X_y,\mathcal O_{X_y})=1$ (相应地, 几何纤维为整概形). 对 $X$ 上的线丛 $L$, 存在唯一的局部闭 (相应地, 闭) 子概形 $Z\subseteq Y$, 使得
\begin{enumerate}[(1)]
\item $X_Z=X\times_Y Z$ 上的 $L_Z$ 为 $Z$ 上某个线丛的拉回;
\item 若 $T\to Y$ 为概形态射, 且 $X_T$ 上的 $L_T$ 为 $T$ 上某个线丛的拉回, 则 $T\to Y$ 经由 $Z$ 分解.
\end{enumerate}
换言之, 函子
\[
\begin{aligned}
\Sch/Y&\longrightarrow\Sets,\\
(T\to Y)&\longmapsto\begin{cases}\{*\},&L_T\text{ 为 }T\text{ 上某个线丛的拉回},\\\varnothing,&\text{否则}
\end{cases}
\end{aligned}
\]
可由 $Y$ 的局部闭 (相应地, 闭) 子概形表出.
\end{proposition}
\begin{proof}
先观察, $L$ 为线丛的拉回, 当且仅当 $f_*L$ 为线丛, 且伴随映射 $f_*f^*L\to L$ 为同构. 事实上, 若 $L=f^*M$, 其中 $M$ 为 $Y$ 上的线丛, 则由投影公式和等式 $\mathcal O_Y=f_*\mathcal O_X$ (引理 \bookref{A.6.13}), 可知
\[
f_*L\simeq f_*f^*M\simeq f_*\mathcal O_X\otimes M\simeq M
\]
为线丛, 且 $f^*f_*L\to L$ 为同构. 同样的结论适用于基变换 $X_T\to T$, 因而问题在 $Y$ 和 $T$ 上是 Zariski 局部的. 我们将证明每个点 $y\in Y$ 都有使命题成立的开邻域.

由半连续性定理 (\bookref{A.6.4}), 点集 $V=\{y\in Y\mid\mathrm h^0(X_y,L_y)\leq1\}$ 开. 由于对 $y\notin V$, $L_y$ 不平凡, 可把 $Y$ 换成 $V$, 并假设对所有 $y\in Y$ 有 $\mathrm h^0(X_y,L_y)\leq1$. 由上同调与基变换 (\bookref{A.6.2}), 把 $Y$ 换成 $y$ 的仿射开邻域后, 可假设存在有限生成自由 $A$-模之间的同态 $d:A^{r_0}\xrightarrow{d}A^{r_1}$, 使得对每个环映射 $A\to B$, 有 $\mathrm H^0(X_B,L_B)=\ker(d\otimes B)$. 用 $d$ 的对偶 $d^\vee$, 定义 $M$ 为序列
\[
A^{r_1}\xrightarrow{d^\vee}A^{r_0}\longrightarrow M\longrightarrow0
\]
中的余核. 沿 $A\to B$ 作张量积得到右正合列
\[
B^{r_1}\xrightarrow{d^\vee\otimes B}B^{r_0}\longrightarrow M\otimes_A B\longrightarrow0;
\]
应用反变左正合函子 $\operatorname{Hom}_B(-,B)$ 后得到
\[
0\longrightarrow\operatorname{Hom}_B(M\otimes_A B,B)\longrightarrow B^{r_0}\xrightarrow{d\otimes_A B}B^{r_1}.
\]
于是
\begin{equation}\label{src:A.6.19}
\mathrm H^0(X_B,L_B)=\operatorname{Hom}_B(M\otimes_A B,B)=\operatorname{Hom}_A(M,B).
\end{equation}
对 $y\in Y$ 的 $A\to\kappa(y)$ 应用此式, 得到 $\mathrm H^0(X_y,L_y)=\operatorname{Hom}_A(M,\kappa(y))=(M\otimes_A\kappa(y))^\vee$.

若 $\mathrm h^0(X_y,L_y)=0$, 则 $L_y$ 不平凡, 且由于 $M\otimes_A\kappa(y)=0$, 存在 $y$ 的开邻域 $U$, 使得 $\widetilde M|_U=0$. 命题在 $U$ 上成立. 若 $\mathrm h^0(X_y,L_y)=1$, 则 $M\otimes_A\kappa(y)=\kappa(y)$; 由 Nakayama 引理, 把 $Y$ 换成 $y$ 的仿射开邻域后, 存在满射 $A\to M$. 写 $M=A/I$, 其中 $I$ 为理想, 并定义闭子概形 $Z=V(I)\subseteq Y$. 注意 $\mathrm H^0(Z,L_Z)=\operatorname{Hom}_A(A/I,A/I)=A/I$, 因而 $f_{Z,*}L_Z$ 为平凡线丛. 为证明 $f_{Z,*}L_Z$ 的构造与基变换交换, 设 $B$ 为 $A/I$-代数, 注意 $\mathrm H^0(X_B,L_B)=\operatorname{Hom}_A(A/I,B)=B$, 且 $\mathrm H^0(X_Z,L_Z)\otimes_{A/I}B\to\mathrm H^0(X_B,L_B)$ 为同构.

断言 $T\to Y$ 经由 $Z$ 分解, 当且仅当 $f_{T,*}L_T$ 为线丛. ($\Rightarrow$) 显然: $f_{Z,*}L_Z$ 为线丛, 且构造与基变换交换. 逆向在 $T$ 上是 Zariski 局部的, 可假设 $T=\Spec B$ 仿射且 $f_{T,*}L_T=\mathcal O_T$. 则 \bookeqref{A.6.19} 意味着 $B=\operatorname{Hom}_A(A/I,B)$. 因而 $I\subseteq\ker(A\to B)$, 即 $A\to B$ 分解为 $A\to A/I\to B$.

最后, 考虑 $X_Z$ 上的伴随态射 $\lambda:f_Z^*f_{Z,*}L_Z\to L_Z$. 断言对 $y\in Z$, $L_y$ 平凡当且仅当 $\lambda|_{X_y}$ 满射. 若 $\lambda|_{X_y}$ 满射, 则用 $f_{Z,*}L_Z=\mathcal O_Z$, 得到线丛之间的满射 $\mathcal O_{X_y}\to L_y$, 因而为同构. 反之, 由于 $f_{Z,*}L_Z$ 与基变换交换, 比较映射 $f_{Z,*}L_Z\otimes\kappa(y)=\mathrm H^0(X_y,L_y)$ 为同构. 记 $f_y:X_y\to\Spec\kappa(y)$, 有等同
\[
(f_Z^*f_{Z,*}L_Z)|_{X_y}=f_y^*(f_{Z,*}L_Z\otimes\kappa(y))=f_y^*f_{y,*}L_y,
\]
而 $\lambda|_{X_y}$ 对应于伴随映射 $f_y^*f_{y,*}L_y\to L_y$, 它是同构 (因为 $L_y$ 平凡). 把 $Z$ 换成 $Z\setminus\operatorname{Supp}(\operatorname{coker}(\lambda))$, 即得命题. 若纤维 $X_y$ 还几何整, 则命题 \bookref{A.6.16} 意味着 $Z$ 闭. 另见 \cite[第 $90$ 页]{Mum70a}, \cite[引理 $1.19$]{Vie95} 和 \cite[Tags 0BEZ, 0BF0]{SP}.
\end{proof}
\begin{remark}\label{src:A.6.20}
对固有平坦态射 $X\to S$, \newterm{相对 Picard 函子}{relative Picard functor} 定义为
\[
\operatorname{Pic}_{X/S}^{\mathrm{rel}}:\Sch/S\longrightarrow\Sets,\qquad T\longmapsto\operatorname{Pic}(X_T)/\operatorname{Pic}(T);
\]
Picard 函子的说明见 \S\bookref{7.4.8-subsection}. 若 $f:X\to S$ 的纤维几何约化 (相应地, 几何整), 则存在由命题 \bookref{A.6.18} 刻画的局部闭 (相应地, 闭) 子概形 $Z\subseteq Y$, 等价于 $\operatorname{Pic}_{X/S}^{\mathrm{rel}}$ 的对角线可由局部闭浸入 (相应地, 闭浸入) 表出. 在几何整纤维的情形下, 这转化为 $\operatorname{Pic}_{X/S}^{\mathrm{rel}}\to S$ 的分离性. 采用这一语言, 上述结果已在 \cite[定理 $3.1$]{FGAV} 中建立.
\end{remark}
